% Generated by publication/build_sections.py; do not edit by hand.
\paragraph{Informal verdict.}
\textbf{A --- Complete solution}, with \textbf{high confidence} for every substantive case \(n\ge 1\) and \(1\le p<n\).

\paragraph{Lean coverage.}
The Lean development proves parts \(a\)--\(d\), the Peano expansion, the differentiation identities, and the order-zero counterexample. It models the interval as the whole real line and documents the resulting endpoint and node-ordering conventions.

\subsection{Audited informal answer}
\subsubsection{1. Faithful restatement and conventions}

Let \(I\subset \mathbb R\) be an interval and \(f:I\to\mathbb R\). For distinct nodes
\[
x_0< x_1<\cdots <x_n
\]
in \(I\), let \(P_{x_0,\ldots,x_n}\) be the Lagrange interpolation polynomial of \(f\) at those nodes, written in the ordinary monomial basis as
\[
P_{x_0,\ldots,x_n}(X)
=\sum_{q=0}^{n}c_q(x_0,\ldots,x_n)X^q.
\]
The function \(f\) is said to admit an interpolated Taylor expansion of order \(n\), abbreviated DIT(\(n\)), at \(x\in I\), when every \(c_q(x_0,\ldots,x_n)\) has a finite limit as all nodes tend to \(x\). The notation \(f^{[q]}(x)\) denotes (q!) times that limit. The four questions are:

\begin{enumerate}
\item Show that it is enough to assume convergence of the highest coefficient \(c_n\).
\item Show that DIT(\(n\)) implies an ordinary asymptotic expansion of order \(n\).
\item If \(p<n\) and \(f\in C^p\), prove that \(f\) has DIT(\(n\)) at \(x\) if and only if \(f^{(p)}\) has DIT(\(n-p\)) there.
\item If \(f\) has DIT(\(p\)) at every point, determine whether DIT(\(n\)) for \(f\) is equivalent to DIT(\(n-p\)) for the resulting function \(f^{[p]}\). 
\end{enumerate}

There are three minor conventions to make explicit.

\begin{itemize}
\item ``Increasing order'' must mean \textbf{strictly} increasing, since ordinary Lagrange interpolation requires distinct nodes.
\item Limits at an endpoint of \(I\) are relative, hence one-sided. Likewise, \(C^p(I)\) means \(C^p\) on the interior, with the corresponding one-sided continuous extensions at finite endpoints.
\item Because the polynomial is expanded in powers of \(X\), rather than \(X-x\), the lower quantities \(f^{[q]}(x)\), \(q<n\), depend on the ambient order \(n\). When necessary, they will be denoted \(f^{[q]}_{(n)}(x)\). The top quantity \(f^{[n]}(x)\) is intrinsic.
\end{itemize}

There is also a genuine order-zero convention issue, addressed at the end.

\medskip\hrule\medskip

\subsubsection{2. Divided differences}

For pairwise distinct \(t_0,\ldots,t_r\in I\), write
\[
[t_0,\ldots,t_r]f
:=
\sum_{i=0}^{r}
\frac{f(t_i)}
{\displaystyle\prod_{\substack{0\le j\le r\\j\ne i}}(t_i-t_j)}.
\]
This is the divided difference of order \(r\). It is symmetric in its arguments.

From the Lagrange formula,
\[
P_{t_0,\ldots,t_r}(X)
=
%
\sum_{i=0}^{r}f(t_i)
\prod_{j\ne i}\frac{X-t_j}{t_i-t_j},
\]
the leading coefficient is
\[
c_r(t_0,\ldots,t_r)=[t_0,\ldots,t_r]f.
\]

We shall repeatedly use the replacement identity
\[
[A,u]f-[A,v]f=(u-v)[A,u,v]f, \tag{2.1}
\]
where \(A\) is any list of distinct common nodes and all displayed nodes are pairwise distinct. It follows immediately from the usual recursion for divided differences, or from the explicit formula above.

\paragraph{Lemma 2.1 --- Descent of divided differences}

Let \(r\ge1\), and suppose that the order-\(r\) divided differences of \(f\) are bounded near \(x\): there exist \(M,\rho_0>0\) such that
\[
|[t_0,\ldots,t_r]f|\le M
\]
whenever the nodes are distinct and belong to
\[
U_{\rho_0}:=I\cap(x-\rho_0,x+\rho_0).
\]
Then the order-(\(r-1\)) divided differences have a finite limit as their \(r\) nodes tend to \(x\).

\subparagraph{Proof}

Fix \(0<\rho<\rho_0\), and let
\[
A=(a_0,\ldots,a_{r-1}),\qquad
B=(b_0,\ldots,b_{r-1})
\]
be two \(r\)-tuples of distinct nodes in \(U_\rho\).

Choose distinct auxiliary nodes
\[
C=(c_0,\ldots,c_{r-1})\subset U_\rho
\]
disjoint from all the nodes occurring in \(A\) and \(B\). This is possible because a nondegenerate real interval contains infinitely many points in every relative neighborhood of \(x\).

Replace the nodes of \(A\) one at a time by the corresponding nodes of \(C\). At each step, (2.1) gives
\[
\left\lvert 
[A^{(j+1)}]f-[A^{(j)}]f
\right\rvert 
\le |c_j-a_j|M
\le 2\rho M.
\]
There are \(r\) replacements, hence
\[
|[A]f-[C]f|\le 2r\rho M.
\]
Similarly,
\[
|[B]f-[C]f|\le 2r\rho M.
\]
Therefore
\[
|[A]f-[B]f|\le 4r\rho M. \tag{2.2}
\]

Thus the oscillation of the order-(\(r-1\)) divided difference over all \(r\)-tuples in \(U_\rho\) tends to zero with \(\rho\). The Cauchy criterion gives a unique finite limit as all nodes tend to \(x\). \(\square\)

Consequently, if the order-\(n\) divided difference has a finite limit at \(x\), then successively all lower divided differences do. We write
\[
\lambda_k(f;x)
:=
\lim_{t_0,\ldots,t_k\to x}[t_0,\ldots,t_k]f,
\qquad 0\le k\le n. \tag{2.3}
\]

For \(n\ge1\), the limit \(\lambda_0(f;x)\) equals \(f(x)\). Indeed,
\[
f(t)-f(x)=(t-x)[x,t]f,
\]
and the first divided differences remain bounded near \(x\).

\medskip\hrule\medskip

\subsection{3. Part \(a\)}

\subsubsection{Theorem 3.1}

For \(n\ge1\), the following are equivalent:

\begin{enumerate}
\item \(f\) has DIT(\(n\)) at \(x\);
\item the highest coefficient \(c_n(x_0,\ldots,x_n)\) has a finite limit as all nodes tend to \(x\);
\item the order-\(n\) divided difference
\end{enumerate}
\[
   [x_0,\ldots,x_n]f
\]
   has a finite limit as all nodes tend to \(x\).

Moreover, if the limits \(\lambda_k(f;x)\) are defined by (2.3), the interpolation polynomials converge coefficientwise to
\[
Q_{n,x}(X)
:=
\sum_{k=0}^{n}\lambda_k(f;x)(X-x)^k. \tag{3.1}
\]

\paragraph{Proof}

The equivalence of 2 and 3 follows from
\[
c_n(x_0,\ldots,x_n)=[x_0,\ldots,x_n]f.
\]

Assume that this highest coefficient has a finite limit. It is then bounded near the diagonal. Lemma 2.1, applied successively, shows that all limits
\[
\lambda_k(f;x),\qquad k=0,\ldots,n,
\]
exist.

For any ordering of the distinct nodes, Newton's interpolation formula is
\[
P_{x_0,\ldots,x_n}(X)
=
%
\sum_{k=0}^{n}
[x_0,\ldots,x_k]f
\prod_{j=0}^{k-1}(X-x_j). \tag{3.2}
\]
As all nodes tend to \(x\),
\[
[x_0,\ldots,x_k]f\longrightarrow \lambda_k(f;x)
\]
and
\[
\prod_{j=0}^{k-1}(X-x_j)\longrightarrow (X-x)^k
\]
coefficientwise. Hence (3.2) converges coefficientwise to \(Q_{n,x}\). In particular, all ordinary monomial coefficients \(c_q\) have finite limits.

The converse is immediate because DIT(\(n\)) includes convergence of \(c_n\). \(\square\)

Expanding (3.1) in powers of \(X\), we obtain the explicit answer
\[
\lim c_q(x_0,\ldots,x_n)
=
%
\sum_{k=q}^{n}
\binom{k}{q}(-x)^{k-q}\lambda_k(f;x). \tag{3.3}
\]
Thus
\[
f^{[q]}_{(n)}(x)
=
%
q!\sum_{k=q}^{n}
\binom{k}{q}(-x)^{k-q}\lambda_k(f;x). \tag{3.4}
\]
In particular,
\[
f^{[n]}(x)=n!\lambda_n(f;x). \tag{3.5}
\]

Formula (3.4) also explains why the lower \(f^{[q]}_{(n)}\) generally depend on \(n\).

\medskip\hrule\medskip

\subsection{4. Part \(b\)}

\subsubsection{Theorem 4.1}

If \(f\) has DIT(\(n\)) at \(x\), with \(n\ge1\), then
\[
f(x+h)
=
%
\sum_{k=0}^{n}\lambda_k(f;x)h^k
+o(|h|^n)
\qquad
(h\to0,\ x+h\in I). \tag{4.1}
\]
Equivalently,
\[
f(t)=Q_{n,x}(t)+o(|t-x|^n), \tag{4.2}
\]
where \(Q_{n,x}\) is the coefficientwise limit of the interpolation polynomials.

\paragraph{Proof}

Fix \(h\ne0\), with \(x+h\in I\). Choose \(n\) pairwise distinct auxiliary nodes
\[
u_0^{(m)},\ldots,u_{n-1}^{(m)}\in I\setminus{x+h}
\]
such that
\[
u_j^{(m)}\longrightarrow x
\qquad(m\to\infty)
\]
for every \(j\).

Apply Newton's formula to the nodes
\[
u_0^{(m)},\ldots,u_{n-1}^{(m)},x+h
\]
and evaluate at \(X=x+h\):
\[
\begin{aligned}
f(x+h)
={}&
\sum_{k=0}^{n-1}
[u_0^{(m)},\ldots,u_k^{(m)}]f
\prod_{j=0}^{k-1}(x+h-u_j^{(m)})\\
&+
[u_0^{(m)},\ldots,u_{n-1}^{(m)},x+h]f
\prod_{j=0}^{n-1}(x+h-u_j^{(m)}).
\end{aligned} \tag{4.3}
\]

As \(m\to\infty\), the terms of order \(k<n\) tend to
\[
\lambda_k(f;x)h^k,
\]
and the final product tends to \(h^n\ne0\). Hence the last divided difference in (4.3) has a limit, namely
\[
R(h):=
\frac{
f(x+h)-\displaystyle\sum_{k=0}^{n-1}\lambda_k(f;x)h^k
}{h^n}. \tag{4.4}
\]

We claim that
\[
R(h)\longrightarrow\lambda_n(f;x).
\]
Given \(\varepsilon>0\), choose \(\delta>0\) such that
\[
\bigl|[t_0,\ldots,t_n]f-\lambda_n(f;x)\bigr|<\varepsilon \tag{4.5}
\]
whenever all the distinct \(t_i\) lie in \(I\cap(x-\delta,x+\delta)\).

If \(0<|h|<\delta\), then, for sufficiently large \(m\), all the auxiliary nodes in (4.3), as well as \(x+h\), lie in this neighborhood. Thus the last divided difference in (4.3) differs from \(\lambda_n(f;x)\) by less than \(\varepsilon\). Passing to its limit gives
\[
|R(h)-\lambda_n(f;x)|\le\varepsilon.
\]
This proves (4.1). \(\square\)

In terms of the coefficients from the statement,
\[
Q_{n,x}(X)
=
%
\sum_{q=0}^{n}\frac{f^{[q]}_{(n)}(x)}{q!}X^q,
\]
so another exact formulation of the requested expansion is
\[
f(t)
=
%
\sum_{q=0}^{n}\frac{f^{[q]}_{(n)}(x)}{q!}t^q
+o(|t-x|^n). \tag{4.6}
\]

The centered coefficients \(\lambda_k(f;x)\) are the more intrinsic ones.

\medskip\hrule\medskip

\subsection{5. Differentiating a DIT}

The central result for parts \(c\) and \(d\) is the following one-step equivalence.

\subsubsection{Lemma 5.1 --- Differentiation theorem}

Let \(r\ge1\), and let \(f\) be \(C^1\) on a relative neighborhood of \(x\). Then
\[
f\text{ has DIT}(r)\text{ at }x
\quad\Longleftrightarrow\quad
f'\text{ has DIT}(r-1)\text{ at }x. \tag{5.1}
\]
When these conditions hold,
\[
\lambda_{r-1}(f';x)=r,\lambda_r(f;x). \tag{5.2}
\]

In fact, continuity of (f') is more than is needed; differentiability on a neighborhood suffices.

\paragraph{Proof: from \(f\) to (f')}

Assume
\[
\lambda_r(f;x)=L.
\]

Let
\[
y_0<\cdots<y_{r-1}
\]
be distinct nodes near \(x\), and define
\[
G(t):=[y_0+t,\ldots,y_{r-1}+t]f.
\]
Because all differences \((y_i+t)-(y_j+t)\) are independent of \(t\), the explicit divided-difference formula gives
\[
G'(0)=[y_0,\ldots,y_{r-1}]f'. \tag{5.3}
\]

For sufficiently small nonzero \(h\), telescope by shifting the nodes one at a time:
\[
\begin{aligned}
G(h)-G(0)
={}&
\sum_{j=0}^{r-1}
\Big(
[y_0+h,\ldots,y_j+h,y_{j+1},\ldots,y_{r-1}]f\\
&\hspace{38mm}
--------------
%
[y_0+h,\ldots,y_{j-1}+h,y_j,\ldots,y_{r-1}]f
\Big).
\end{aligned}
\]
By the replacement identity (2.1), each summand equals
\[
h[Z_j(h)]f,
\]
where \(Z_j(h)\) is the (\(r+1\))-tuple consisting of the common \(r-1\) nodes together with \(y_j\) and \(y_j+h\). Hence
\[
\frac{G(h)-G(0)}h
=
%
\sum_{j=0}^{r-1}[Z_j(h)]f. \tag{5.4}
\]

Let \(\varepsilon>0\). If all \(y_i\) are sufficiently close to \(x\), then for all sufficiently small \(h\), every divided difference on the right of (5.4) differs from \(L\) by at most \(\varepsilon/r\). Therefore
\[
\left\lvert 
\frac{G(h)-G(0)}h-rL
\right\rvert <\varepsilon.
\]
Letting \(h\to0\) and using (5.3),
\[
\bigl|[y_0,\ldots,y_{r-1}]f'-rL\bigr|
\le\varepsilon.
\]
Thus
\[
[y_0,\ldots,y_{r-1}]f'\longrightarrow rL,
\]
which, by part \(a\), is precisely DIT(\(r-1\)) for (f').

\paragraph{Proof: from (f') to \(f\)}

Conversely, assume
\[
\lambda_{r-1}(f';x)=M.
\]

Take arbitrary nodes
\[
z_0<z_1<\cdots<z_r
\]
near \(x\), and let \(P\) be the degree-\(\le r\) interpolation polynomial of \(f\) at these nodes. The function
\[
H:=f-P
\]
vanishes at each \(z_i\). By Rolle's theorem, for every \(i=0,\ldots,r-1\), there is a point
\[
\xi_i\in(z_i,z_{i+1})
\]
such that
\[
H'(\xi_i)=0.
\]
Thus
\[
P'(\xi_i)=f'(\xi_i).
\]
The points satisfy
\[
\xi_0<\cdots<\xi_{r-1},
\]
and (P'), of degree at most \(r-1\), is therefore the interpolation polynomial of (f') at these \(r\) points.

The leading coefficient of (P') is
\[
r[z_0,\ldots,z_r]f,
\]
while the leading coefficient of the interpolation polynomial of (f') is
\[
[\xi_0,\ldots,\xi_{r-1}]f'.
\]
Consequently,
\[
r[z_0,\ldots,z_r]f
=
%
[\xi_0,\ldots,\xi_{r-1}]f'. \tag{5.5}
\]

As all \(z_i\to x\), all \(\xi_i\to x\). Therefore the right-hand side of (5.5) tends to \(M\), and
\[
[z_0,\ldots,z_r]f\longrightarrow \frac Mr.
\]
Part \(a\) now gives DIT(\(r\)) for \(f\), with
\[
\lambda_r(f;x)=\frac Mr.
\]
This completes the proof. \(\square\)

\medskip\hrule\medskip

\subsection{6. Part \(c\)}

Let \(0\le p<n\) and \(f\in C^p(I)\).

For \(p=0\), the assertion is tautological. Suppose \(p\ge1\). Apply Lemma 5.1 successively to
\[
f,\ f',\ f'',\ldots,f^{(p-1)}
\]
with orders
\[
n,\ n-1,\ n-2,\ldots,n-p+1.
\]
We obtain
\[
\begin{aligned}
f\text{ has DIT}(n)\text{ at }x
&\Longleftrightarrow f'\text{ has DIT}(n-1)\text{ at }x\\
&\Longleftrightarrow f''\text{ has DIT}(n-2)\text{ at }x\\
&\quad\vdots\\
&\Longleftrightarrow f^{(p)}\text{ has DIT}(n-p)\text{ at }x.
\end{aligned}
\]

This proves the required equivalence.

The limiting highest coefficients satisfy the precise identity
\[
\lambda_{n-p}(f^{(p)};x)
=
%
= n(n-1)\cdots(n-p+1)\lambda_n(f;x)
%
\frac{n!}{(n-p)!}\lambda_n(f;x). \tag{6.1}
\]
Multiplying by (\(n-p\)!), we get the particularly simple relation
\[
\boxed{
\bigl(f^{(p)}\bigr)^{[,n-p,]}(x)
=
%
f^{[n]}(x).
} \tag{6.2}
\]

Thus part \(c\) holds, and (6.1) gives the exact transformation of the top divided-difference limit.

\medskip\hrule\medskip

\subsection{7. Part \(d\)}

The substantive answer is affirmative for every \(1\le p<n\).

The point is that having DIT(\(p\)) at every point automatically forces \(f\) to be \(C^p\), and \(f^{[p]}\) is then exactly the ordinary derivative \(f^{(p)}\).

\subsubsection{Lemma 7.1 --- DIT(\(1\)) everywhere implies \(C^1\)}

Suppose \(g\) has DIT(\(1\)) at every point of \(I\). Then \(g\in C^1(I)\), and
\[
g'(x)=\lambda_1(g;x)=g^{[1]}(x).
\]

\paragraph{Proof}

At a fixed \(x\),
\[
\frac{g(t)-g(x)}{t-x}=[x,t]g\longrightarrow\lambda_1(g;x),
\]
so \(g\) is differentiable and
\[
g'(x)=\lambda_1(g;x).
\]

It remains to prove continuity of (g'). Given \(\varepsilon>0\), the strict secant limit at \(x\) provides \(\delta>0\) such that
\[
|[u,v]g-g'(x)|<\varepsilon
\]
whenever \(u\ne v\) and both (u,v) lie within \(\delta\) of \(x\).

For \(y\) sufficiently close to \(x\), let \(v\to y\). All secants ([y,v]g) are then within \(\varepsilon\) of (g'\(x\)). Passing to the limit \(v\to y\) gives
\[
|g'(y)-g'(x)|\le\varepsilon.
\]
Thus (g') is continuous. \(\square\)

\subsubsection{Proposition 7.2}

Let \(p\ge1\). If \(f\) has DIT(\(p\)) at every point of \(I\), then
\[
f\in C^p(I)
\]
and
\[
\boxed{f^{[p]}=f^{(p)}.} \tag{7.1}
\]

\paragraph{Proof}

Since DIT(\(p\)) implies, by the descent argument in part \(a\), existence of the first divided-difference limit, \(f\) has DIT(\(1\)) at every point. Lemma 7.1 therefore gives
\[
f\in C^1(I).
\]

By Lemma 5.1, because \(f\) has DIT(\(p\)) everywhere,
\[
f'
\]
has DIT(\(p-1\)) everywhere. If \(p-1\ge1\), Lemma 7.1 applied to (f') gives \(f'\in C^1\), hence \(f\in C^2\).

Continuing inductively,
\[
f^{(j)}
\]
has DIT(\(p-j\)) everywhere and is \(C^1\) for \(j=0,\ldots,p-1\). Hence
\[
f\in C^p(I).
\]

The factor relation (5.2), iterated \(p\) times, gives
\[
f^{(p)}(x)=p!\lambda_p(f;x).
\]
Since \(c_p\) is the order-\(p\) divided difference,
\[
f^{[p]}(x)=p!\lambda_p(f;x).
\]
Therefore
\[
f^{[p]}(x)=f^{(p)}(x)
\]
at every point. \(\square\)

We may now apply part \(c\):

\[
\begin{aligned}
f\text{ has DIT}(n)\text{ at }x
&\Longleftrightarrow
f^{(p)}\text{ has DIT}(n-p)\text{ at }x\\
&\Longleftrightarrow
f^{[p]}\text{ has DIT}(n-p)\text{ at }x.
\end{aligned}
\]

Hence:

\[
\boxed{
\text{For every }1\le p<n,\text{ the answer to part (d) is yes.}
} \tag{7.2}
\]

The global assumption ``at every point of \(I\)'' can be weakened: for the equivalence at a fixed \(x\), it is enough that \(f\) have DIT(\(p\)) throughout a relative neighborhood of \(x\).

\medskip\hrule\medskip

\subsubsection{8. The order-zero ambiguity}

If \(p=0\), the answer depends on what ``DIT(\(0\))'' means.

For \(n=0\), the interpolation polynomial is simply
\[
P_{x_0}(X)=f(x_0).
\]
Thus the literal condition in the PDF is merely the existence of
\[
\lim_{t\to x}f(t).
\]

\paragraph{Continuity convention}

If DIT(\(0\)) is understood to require
\[
\lim_{t\to x}f(t)=f(x),
\]
then
\[
f^{[0]}=f,
\]
and part \(d\) for \(p=0\) is tautologically true.

\paragraph{Punctured-limit convention}

If only the finite punctured limit is required, part \(d\) is false for \(p=0\).

Take
\[
I=(-1,1),\qquad
f(t)=
\begin{cases}
1,&t=0,\\
0,&t\ne0.
\end{cases}
\]
At every point, the punctured limit exists, and
\[
f^{[0]}\equiv0.
\]
Thus \(f^{[0]}\) has DIT(\(1\)) at \(0\). But \(f\) does not, since for \(t>0\),
\[
[0,t]f=\frac{f(t)-f(0)}{t}=-\frac1t,
\]
which diverges.

Therefore the fully literal conclusion is:

\[
\boxed{
\begin{array}{l}
\text{Part (d) is affirmative for every }p\ge1.[2mm]
\text{For }p=0,\text{ it is affirmative under the continuity convention,}\\
\text{but false under the punctured-limit convention.}
\end{array}}
\]

The same convention affects part \(b\) when \(n=0\): a punctured finite limit gives an expansion with a constant equal to that limit, but not necessarily with constant \(f(x)\).

\medskip\hrule\medskip

\subsection{9. Checks and sharpness}

\paragraph{9.1 The case \(n=1\)}

DIT(\(1\)) at \(x\) means
\[
[u,v]f=\frac{f(v)-f(u)}{v-u}
\]
has a limit as \(u,v\to x\). This is strict differentiability at \(x\). The results give
\[
f(x+h)=f(x)+\lambda_1(f;x)h+o(|h|),
\]
and, if this condition holds everywhere, \(f\in C^1\).

\paragraph{9.2 Polynomial check}

For a polynomial \(f\), all the limits exist. The limiting interpolation polynomial is its Taylor polynomial of degree \(n\) at \(x\):
\[
Q_{n,x}(X)
=
%
\sum_{k=0}^{n}\frac{f^{(k)}(x)}{k!}(X-x)^k,
\]
with the sum truncated automatically if \(\deg f<n\). Thus
\[
\lambda_k(f;x)=\frac{f^{(k)}(x)}{k!},
\]
and (6.1) becomes
\[
\frac{f^{(n)}(x)}{(n-p)!}
=
%
\frac{n!}{(n-p)!}\frac{f^{(n)}(x)}{n!},
\]
as required.

\paragraph{9.3 The converse to part \(b\) is false}

An ordinary asymptotic expansion need not imply a DIT.

Let
\[
g(0)=0,\qquad
g(t)=t^2\sin(1/t)\quad(t\ne0).
\]
Then
\[
g(t)=o(t),
\]
so \(g\) has a first-order expansion at \(0\) with derivative \(0\). But
\[
g'(t)=2t\sin(1/t)-\cos(1/t)
\]
oscillates between values close to \(1\) and (-1) as \(t\to0\).

At points \(t_k\to0\) where \(g'(t_k)\to1\), choose \(h_k\to0\) so that
\[
[t_k,t_k+h_k]g\to1.
\]
At another sequence \(s_k\to0\), choose secants tending to (-1). Hence the two-point divided differences do not have a unique limit as both nodes tend to \(0\). Thus \(g\) does not have DIT(\(1\)) there.

So DIT is strictly stronger than the existence of an ordinary Peano--Taylor expansion.

\paragraph{9.4 Role of the hypotheses}

\begin{itemize}
\item Distinctness of the nodes is indispensable for ordinary Lagrange interpolation. Repeated nodes would require Hermite interpolation.
\item The interval assumption is used essentially in the Rolle-theorem argument: each consecutive pair of interpolation nodes determines an intervening point.
\item The \(C^p\) assumption in part \(c\) is sufficient but not optimal. The proof only needs \(p\)-fold differentiability on a neighborhood of \(x\).
\item In part \(d\), the ``every point'' hypothesis is what upgrades the strict divided-difference regularity to \(C^p\). Locally around \(x\) is sufficient.
\item Endpoint cases require no separate argument beyond taking relative limits and one-sided derivatives.
\end{itemize}

\medskip\hrule\medskip

\subsection{Final classification}

For \(n\ge1\), let
\[
\Lambda_n(f;x)
=
%
\lim_{x_0,\ldots,x_n\to x}[x_0,\ldots,x_n]f.
\]

Then:

\[
\boxed{
f\text{ has DIT}(n)\text{ at }x
\iff
\Lambda_n(f;x)\in\mathbb R.
}
\]

When this holds, all lower limits \(\lambda_k(f;x)\) exist, the interpolation polynomials converge to
\[
Q_{n,x}(X)=\sum_{k=0}^{n}\lambda_k(f;x)(X-x)^k,
\]
and
\[
\boxed{
f(x+h)=\sum_{k=0}^{n}\lambda_k(f;x)h^k+o(|h|^n).
}
\]

For \(0\le p<n\) and \(f\in C^p\),
\[
\boxed{
f\text{ has DIT}(n)\text{ at }x
\iff
f^{(p)}\text{ has DIT}(n-p)\text{ at }x,
}
\]
with
\[
\boxed{
\lambda_{n-p}(f^{(p)};x)
=
%
\frac{n!}{(n-p)!}\lambda_n(f;x).
}
\]

Finally, if \(p\ge1\) and \(f\) has DIT(\(p\)) at every point, then
\[
\boxed{
f\in C^p,\qquad f^{[p]}=f^{(p)},
}
\]
and therefore
\[
\boxed{
f\text{ has DIT}(n)\text{ at }x
\iff
f^{[p]}\text{ has DIT}(n-p)\text{ at }x.
}
\]
